\begin{proposition}[2.8.4]
\label{IV.2.8.4-fr}
Bajo las hipótesis sobre $Y$ de \textup{(2.8.1)}, sean $X_1$, $X_2$ dos
$Y$-preesquemas, $\mathcal F_i$ un $\mathcal O_{X_i}$-módulo
cuasicoherente, $\mathcal G'_i$ un $\mathcal O_{(X_i)_\eta}$-módulo cociente
de $(\mathcal F_i)_\eta$ $(i=1,2)$. Entonces se tiene
\end{proposition}

\begin{env}[2.8.4.1]
\label{IV.2.8.4.1-fr}
\[
\overline{\mathcal G'_1}\otimes_{\mathcal O_Y}\overline{\mathcal G'_2}
=\overline{\mathcal G'_1\otimes_{\mathbf{k}(\eta)}\mathcal G'_2}.
\]
\end{env}

\begin{proof}
En efecto, pongamos $X=X_1\times_Y X_2$; el primer miembro de (2.8.3.1) es
un $\mathcal O_X$-módulo cuasicoherente que es $Y$-\emph{plano}
($\mathbf{0}_{\mathrm I}$, 6.2.1), cuya imagen inversa en $X_\eta$ es
$\mathcal G'_1\otimes_{\mathbf{k}(\eta)}\mathcal G'_2$ (\textbf I, 9.1.5) y que es un módulo
cociente de $\mathcal F_1\otimes_{\mathcal O_Y}\mathcal F_2$; por tanto, la
conclusión se sigue de la propiedad de unicidad de (2.8.1).
\end{proof}

